\scanpage{10}
obtenue à partir de l'extension (14) de $\mathbb Z/2\mathbb Z\simeq\mu_2$ par $\mathbb Z$

\diagram{\begin{tikzcd}[ampersand replacement=\&, row sep=small]
0\arrow[r,no head]\&\mathbb Z\arrow[r,no head,"2"]\&\mathbb Z\arrow[r]\&\mathbb Z/2\mathbb Z\arrow[r,no head]\arrow[d,phantom,"\Vert"]\&0\\
\&\&\&\mu_2\&
\end{tikzcd}}

en prenant l'image inverse de celle-ci par l'hom signature
\[
\mathfrak S_I\xrightarrow{\operatorname{sgn}}\mu_2\simeq\mathbb Z/2\mathbb Z,
\]
et en prenant la $\mathbb Z^I$-extension déduite de celle-ci par \unsure{l'hom} \inserted{diag} $\mathbb Z\;-\;\mathbb Z^I$.

Cette extension n'est pas scindable -- car l'extension induite sur un sous-groupe $\gamma_k$ d'ordre $2$, défini par une transposition $\sigma_k=(i,j)$, ne l'est pas -- car en regardant la projection $\mathbb Z^I\;-\;\mathbb Z$ correspondant au troisième indice $k$ (fixe par l'involution) on trouve une extension de $\gamma_k$ par $\mathbb Z$, qui n'est autre (via l'identification $\gamma_k\simeq\mathbb Z/2\mathbb Z$) que l'extension (14) de $\mathbb Z/2\mathbb Z$ par $\mathbb Z$, qui n'est pas scindable évidemment. Par contre la restriction de l'extension à $\mathfrak S_I^+$, qui est le stabilisateur de $x$ dans $\mathfrak S_I$, est évidemment canoniquement scindée -- ce qui est manifeste aussi sur la forme explicite (17)(18) -- le scindage est
\[
(19)\qquad\alpha\longmapsto(0,\alpha)\qquad\alpha\in\mathfrak S_I^+.
\]
\mspage{9}
